Un conducteur $MN$, parcouru par un courant d'intensité $I$ est placé
dans un champ magnétique $\vec{B}$ uniforme.

Dans chacun des cas suivants, \textbf{représenter} la force
électromagnétique qui s'exerce sur le conducteur et
\textbf{calculer} l'intensité de cette force.
\textbf{Données~:}
$I=\qty{1}{\A}$~; $MN=\qty{10}{\cm}$~;
$B=\qty{1e-4}{\T}$.~\textbf{(4pt)}


\begin{tcbitemize}[raster columns=4,raster rows=1,
        raster equal height=rows,raster force size=false,
        raster every box/.style={top=1mm,bottom=0mm}]
    \tcbitem
        \textbf{a)}
        \begin{figure}[H]
            \centering
            \begin{tikzpicture}[thick]
                \node[fil=3cm,anchor=south] (fil) at (0,0) {};
                \node[below] at (fil.south) {$N$};
                \node[above] at (fil.north) {$M$};
                \draw[->,red] (fil.center) --
                    node[left] {$I$} ++(0,.5);
                \node[circle,inner sep=1mm,draw,line width=1.2pt,
                      below right=7mm,
                      label=below:$\vec{B}$] (c) at (fil.east) {};
                \node[circle,inner sep=1mm,fill,
                      inner sep=1pt] at (c.center) {};
            \end{tikzpicture}
        \end{figure}
    \tcbitem
        \textbf{b)}
        \begin{figure}[H]
            \centering
            \begin{tikzpicture}[thick]
                \node[fil=3cm,anchor=south] (fil) at (0,0) {};
                \node[below] at (fil.south) {$N$};
                \node[above] at (fil.north) {$M$};
                \draw[->,red] (fil.center) --
                    node[left] {$I$} ++(0,.5);
                \draw[->] ([xshift=5mm,yshift=-5mm]fil.center) --
                    node[right] {$\vec{B}$} ++(0,1);
            \end{tikzpicture}
        \end{figure}
    \tcbitem
        \textbf{c)}
        \begin{figure}[H]
            \centering
            \begin{tikzpicture}[thick]
                \node[fil=3cm,anchor=south] (fil) at (0,0) {};
                \node[below] at (fil.south) {$N$};
                \node[above] at (fil.north) {$M$};
                \draw[->,red] (fil.center) --
                    node[left] {$I$} ++(0,.5);
                \node[circle,inner sep=1mm,draw,line width=1.2pt,
                      below right=7mm,
                      label=below:$\vec{B}$] (c) at (fil.east) {};
                \draw[line width=1.2pt]
                    (c.45) -- (c.225) (c.135) -- (c.315);
            \end{tikzpicture}
        \end{figure}
    \tcbitem
        \textbf{d)}
        \begin{figure}[H]
            \centering
            \begin{tikzpicture}[thick]
                \node[fil=3cm,anchor=south] (fil) at (0,0) {};
                \node[below] at (fil.south) {$N$};
                \node[above] at (fil.north) {$M$};
                \draw[->,red] (fil.center) --
                    node[left] {$I$} ++(0,.5);
                \draw[->] ([xshift=2mm,yshift=-7mm]fil.center)
                  coordinate (tmp) --
                  node[right=2mm] {$\vec{B}$} ++(45:1.5);
                \draw[shorten <=1pt] ([shift=(45:1)]tmp) arc (45:90:1)
                  node[pos=0.3,above=1mm] {$\ang{45}$};;
            \end{tikzpicture}
        \end{figure}
\end{tcbitemize}

